%% Packages and macros shared by both class wrappers.
%%
%% Nothing class-specific belongs here. The wrappers load their own document
%% class first, then \input this file, then bind the \wtable / \ntable float
%% environments to whatever suits their column layout.

%% Fonts are deliberately NOT set here. Tectonic is XeTeX, so UTF-8 input needs
%% no inputenc, and the two classes want different faces: elsarticle preprints
%% are conventionally Latin Modern, while IEEE requires Times. Each wrapper
%% makes its own choice.

\usepackage{booktabs}
\usepackage{array}
\usepackage{makecell}   % stacked table headers
\usepackage{amsmath}
\usepackage{xcolor}
\usepackage{url}
\usepackage{microtype}  % better line breaking; fewer overfull lines
\usepackage{adjustbox}
\usepackage{tikz}
\usetikzlibrary{arrows.meta, positioning, fit, backgrounds, calc}

% Several results tables carry more numeric columns than the measure holds.
% \fittable shrinks only those that actually overflow and leaves the rest at
% full size, so no table is dropped or relegated to an appendix to make it fit
% --- every number stays on the page it argues from. Measuring against
% \linewidth rather than \textwidth is what makes this work in both layouts:
% inside a full-width float it resolves to the text width, inside a
% single-column float to the column width.
% The scale cap is what makes the tables coherent. Without it adjustbox shrinks
% each table by whatever factor that table alone needs, so a wide one renders at
% 6.1pt beside a narrow one at 8pt. Scaling everything by 0.85 first and only
% then fitting to the column means every table lands at one size, and the widest
% table is what sets that size.
\newenvironment{fittable}
  {\begin{adjustbox}{scale=0.85, max width=\linewidth}}
  {\end{adjustbox}}

% Numeric tables are dense; buy back horizontal room before any scaling.
\setlength{\tabcolsep}{3.5pt}

% Stacked headers are centred; the default \makecell alignment fights the
% right-aligned numeric columns underneath, so pin it explicitly.
\renewcommand{\cellalign}{cc}

% Protocol field names and token literals appear constantly; give them one
% consistent treatment so evidence is always visually distinct from prose.
% Note the deliberate absence of hyphenation inside \code: this paper's
% argument turns on exact token spellings, and a break-hyphen inserted into
% mqtt.protoname would be indistinguishable from part of the identifier.
%% Courier runs wide beside Times, so code spans are set a little smaller than
%% their surroundings. \small is an absolute size, which inside a \footnotesize
%% table renders code LARGER than the table text and inflates the column it sits
%% in; scaling off \f@size keeps the reduction relative wherever it appears.
\makeatletter
\newcommand{\code}[1]{{\ttfamily
  \@tempdima=\f@size pt \@tempdima=0.92\@tempdima
  \fontsize{\@tempdima}{1.2\@tempdima}\selectfont #1}}
\makeatother

%% IEEEtran sets the journal abstract in \bfseries, but math mode does not
%% inherit the font series, so a value written as $1.0000 \pm 0.0000$ prints in
%% the regular weight beside bold text. These two set the symbols in text mode
%% instead, so a mean-and-deviation reads in whatever weight surrounds it under
%% both classes.
\newcommand{\pmv}[2]{#1\,\textpm\,#2}
\newcommand{\by}{\texttimes}

% Long dotted identifiers are unbreakable words, so a few lines cannot be set
% at natural spacing. Let TeX loosen interword space as a last resort rather
% than overflowing the measure.
\emergencystretch=3em
